\documentclass[11pt]{article}
\usepackage[T1]{fontenc}
\usepackage{amsmath,amssymb,amsthm}
\usepackage[margin=1in]{geometry}
\usepackage{hyperref}
\newtheorem{theorem}{Theorem}
\newtheorem{lemma}[theorem]{Lemma}
\newtheorem{proposition}[theorem]{Proposition}
\newtheorem{example}[theorem]{Example}
\newcommand{\dp}{d'}
\newcommand{\ddp}{d''}
\newcommand{\D}{d'd''}
\newcommand{\BT}{\mathrm{BT}}
\title{Closed smooth multipliers of Bedford--Taylor currents\newline
and an explicit analytic test of the closedness condition}
\author{Research working paper for Problem P21.3}
\date{8 October 2026; continuation 367; independent review pending}
\begin{document}
\maketitle
\begin{abstract}
Let $T$ be a positive, doubly closed ordinary Chambert--Loir--Ducros
current which satisfies potential independence for a prescribed class of
continuous potentials. For a smooth nonnegative function $a$, we prove
that $aT$ satisfies the same property if and only if it is doubly closed.
The proof establishes the distributional identity
$\D(a hT)=a\D(hT)$ for every continuous $h$, under precisely the
closedness hypothesis on $aT$. We also treat local uniform limits of such
closed multipliers of a fixed current. On the analytic two-dimensional
torus, a compactly supported test gives an explicit failure of closedness
for an arbitrary positive smooth multiplier of the integration current.
These results are conditional stability statements and a source-scope
diagnostic. They do not solve the assertion that every positive doubly
closed current is Bedford--Taylor.
\end{abstract}

\section{Ordinary currents and the scope of potential independence}
Let $K$ be a complete non-Archimedean field, with trivial valuation
permitted in the general statements. Let $X$ be a good, separated,
boundaryless, pure $n$-dimensional $K$-analytic space. All currents and
test forms below are ordinary CLD currents and smooth forms on ordinary
open subsets. Write $D=\dp\ddp$. A current of bidegree $(p,p)$ acts on
compactly supported smooth forms of type $(n-p,n-p)$. Doubly closed
means $\dp T=\ddp T=0$.

For each ordinary open $U$, fix a class $\mathcal F(U)\subset C(U,\mathbb R)$
which is compatible with restriction. A positive doubly closed current
$T$ has property $\BT_{\mathcal F}$ if, on every such $U$, every pair
$u,v\in\mathcal F(U)$ with $D[u]=D[v]$ satisfies
\[
 D(uT|_U)=D(vT|_U).
\]
Here $[u]=u\delta_U$ is the canonical function current. This formulation
can use all finite continuous individually psh functions in the original
P21.3 problem. It can also use the locally uniformly generated continuous
psh class of CLD. No identification of those classes is made or needed.
The conclusions hold for any prescribed class as above.

Multiplication of a positive current by a continuous function is the
ordinary measure-coefficient operation of CLD, Sections 7.2.11--7.2.12.
It remains defined for a signed locally bounded continuous function by
taking its positive and negative parts. For a smooth multiplier it agrees
with the dual operation on test forms. Smooth-form multiplication of a
current is defined by duality on compactly supported smooth tests.

\section{A distributional multiplication calculation}
\begin{lemma}\label{lem:coeff}
Let $T$ be positive. If $h$ is continuous and $\beta$ is a smooth form,
then for every compact smooth test $\omega$ one has
$\langle\beta\wedge(hT),\omega\rangle
=\int h\,d\mu_{T,\beta\wedge\omega}$, with the usual wedge convention.
In particular, $a(hT)=(ah)T$ for smooth $a$. If
$\beta\wedge T=0$, then $\beta\wedge(hT)=0$. For a net $h_i$ converging
locally uniformly to $h$, one has $h_iT\to hT$ weakly.
\end{lemma}
\begin{proof}
On the compact support of a smooth test $\theta$, a finite smooth
partition of unity reduces to finitely many tropical charts. On each
chart, the symmetric part of $\theta$ is a difference of two positive
forms with compact support: add and subtract a sufficiently large
multiple of a strongly positive standard form times a compact cutoff.
The required domination estimate is the one in CLD Proposition 7.2.9.
The antisymmetric part evaluates to zero on the symmetric positive $T$.
CLD 7.2.11 therefore gives a finite signed Radon measure $\mu_{T\theta}$
such that
$\int\phi\,d\mu_{T\theta}=\langle T,\phi\theta\rangle$
for every smooth scalar $\phi$. The measure is supported on the compact
support of $\theta$. Uniqueness follows because compact smooth scalars
are dense in compact continuous scalars, by CLD Corollary 6.2.6.
The continuous multiplication construction then gives
$\langle hT,\theta\rangle=\int h\,d\mu_{T\theta}$.

Apply this to $\theta=\beta\wedge\omega$ to obtain the displayed
formula. If $\beta\wedge T=0$, evaluation on every $\phi\omega$ shows
$\int\phi\,d\mu_{T,\beta\wedge\omega}=0$. The same uniqueness gives
$\mu_{T,\beta\wedge\omega}=0$, so the displayed formula vanishes for
every $\omega$. This proves the zero assertion without assuming that an
arbitrary order-zero current admits arbitrary continuous multiplication.
For a smooth scalar $a$, uniqueness also gives
$\mu_{T,a\omega}=a\mu_{T\omega}$. Consequently
$\langle a(hT),\omega\rangle=\int h\,d\mu_{T,a\omega}
=\int ah\,d\mu_{T\omega}=\langle(ah)T,\omega\rangle$.

For a fixed test $\omega$ there is therefore a finite signed measure
$\mu_{T\omega}$ supported on its compact support $C$, and
\[
 |\langle(h_i-h)T,\omega\rangle|
 \leq \sup_C|h_i-h|\,|\mu_{T\omega}|(C)\longrightarrow0.
\]
This proves the net assertion without extracting a sequence or making an
assumption about a changing family of currents.
\end{proof}

\begin{lemma}\label{lem:leibniz}
For any current $S$ of total degree $m$ and smooth form $\beta$,
ordinary current differentiation obeys
\[
 \dp(\beta\wedge S)=\dp\beta\wedge S+
                 (-1)^{\deg\beta}\beta\wedge\dp S,
\]
and the same formula with $\ddp$. In particular, for a smooth scalar $a$,
\begin{equation}\label{eq:second}
 D(aS)=Da\wedge S-\ddp a\wedge\dp S
                    +\dp a\wedge\ddp S+aDS.
\end{equation}
These statements do not require $S$ to be closed.
\end{lemma}
\begin{proof}
The right-wedge version is CLD Lemma 7.7.7, whose proof applies the smooth
Leibniz rule to $\beta\wedge\omega$ for a compact smooth test $\omega$.
Moving $\beta$ across $S$ with the total-degree sign gives the displayed
left-wedge version. The $\ddp$ version follows by the same test calculation,
or by the $J$-conjugation identity in CLD 7.7.3. The anticommutation
$\ddp\dp=-\dp\ddp$ holds for all currents, not just closed ones.
Apply $\ddp(aS)=\ddp a\wedge S+a\ddp S$, then apply $\dp$ and its
graded Leibniz rule. The degree of $\ddp a$ is one, giving exactly
\eqref{eq:second}. All products on the right have smooth form factors;
no derivative of the continuous coefficient of an order-zero current is
multiplied by another nonsmooth coefficient.
\end{proof}

\begin{proposition}\label{prop:identity}
Let $T$ be a positive doubly closed current, and let $a$ be smooth such
that $aT$ is doubly closed. For every continuous function $h$,
\begin{equation}\label{eq:main}
 D((ah)T)=aD(hT).
\end{equation}
\end{proposition}
\begin{proof}
Since $T$ is doubly closed, the first-order Leibniz rules give
\[
 \dp a\wedge T=\dp(aT)=0,
 \qquad \ddp a\wedge T=\ddp(aT)=0.
\]
Differentiating the second zero identity yields
\[
 0=\dp(\ddp a\wedge T)=Da\wedge T-\ddp a\wedge\dp T
   =Da\wedge T.
\]
Set $S=hT$. The coefficient lemma, which uses only the positivity and
order-zero multiplication of the fixed $T$, now gives all three identities
\[
 \dp a\wedge S=0,\qquad \ddp a\wedge S=0,
 \qquad Da\wedge S=0.
\]
Even though $S$ need not be closed, its ordinary current differentials
exist. Differentiate its first zero identity by $\ddp$:
\[
 0=\ddp(\dp a\wedge S)
  =-Da\wedge S-\dp a\wedge\ddp S.
\]
Thus $\dp a\wedge\ddp S=0$. Differentiate its second zero identity
by $\dp$:
\[
 0=\dp(\ddp a\wedge S)
  =Da\wedge S-\ddp a\wedge\dp S.
\]
Thus $\ddp a\wedge\dp S=0$ as well. Every extra term on the right of
\eqref{eq:second} vanishes. Hence $D(aS)=aDS$. The coefficient lemma
identifies $aS$ with $(ah)T$, proving \eqref{eq:main}. This derivation
uses no regularization of $h$, no symmetry formula involving two
nonsmooth factors, and no positivity assertion for $D(hT)$.
\end{proof}

\section{The exact smooth multiplier criterion}
\begin{theorem}\label{thm:criterion}
Suppose $T$ has property $\BT_{\mathcal F}$ and $a\geq0$ is smooth.
Then $aT$ has property $\BT_{\mathcal F}$ if and only if $aT$ is doubly
closed. Equivalently, the allowable smooth nonnegative multipliers are
exactly those satisfying
\[
 \dp a\wedge T=\ddp a\wedge T=0.
\]
This is a local criterion on ordinary open subsets and holds over all
complete $K$, including trivially valued $K$.
\end{theorem}
\begin{proof}
Necessity is part of the definition of $\BT_{\mathcal F}$. The displayed
equivalence with double closedness follows from the first-order Leibniz
rules and double closedness of $T$.

For sufficiency, $aT$ is positive by nonnegativity of $a$ and is doubly
closed by assumption. Fix an arbitrary ordinary open $U$ and an arbitrary
admissible pair $u,v\in\mathcal F(U)$ with $D[u]=D[v]$. The hypothesis
on $T$ gives $D((u-v)T|_U)=0$. Apply Proposition \ref{prop:identity}
on $U$ with $h=u-v$. The result is
\[
 D((u-v)aT|_U)=aD((u-v)T|_U)=0,
\]
where associativity is supplied by Lemma \ref{lem:coeff}. This is the
required equality for $aT$ on every ordinary open. No approximation or
generating hypothesis for the admissible pair was used.
\end{proof}

\begin{theorem}\label{thm:limits}
Let $T$ have property $\BT_{\mathcal F}$, and let $a\geq0$ be continuous.
Suppose locally on every point there is a net of smooth nonnegative
functions $a_i\to a$ locally uniformly, with every $a_iT$ doubly closed
on that neighborhood. Then $aT$ has property $\BT_{\mathcal F}$.
\end{theorem}
\begin{proof}
Work in one of the specified ordinary neighborhoods. The preceding
theorem gives $\BT_{\mathcal F}$ for each $a_iT$. For a compact smooth
test $\omega$, Lemma \ref{lem:coeff} gives $a_iT\to aT$ weakly.
First-order differential operators are weakly continuous by their
definition on smooth tests; their limits show $\dp(aT)=\ddp(aT)=0$.
Nonnegative multiplication shows $aT$ is positive.

For an admissible pair $u,v$ on any ordinary open subset of this
neighborhood, fix a compact test $\omega$. The supports of $\omega$ and
$D\omega$ are compact. With $h=u-v$, associativity and the fixed-current
measure formula give
\[
 \langle D((a_i-a)hT),\omega\rangle
 =\int (a_i-a)h\,d\mu_{T,D\omega}\longrightarrow0.
\]
Indeed $h$ is bounded on that compact support and $a_i\to a$ uniformly
there. Each $D(a_i hT)$ is zero, so $D(ahT)=0$. This proves the property
locally; the zero equality glues as an equality of ordinary currents.
For completeness, a compact test on an arbitrary open can be split by a
finite smooth partition of unity subordinate to these neighborhoods.
Every local zero current then evaluates to zero on each summand. Positivity
and double closedness are also local. Thus $aT$ has the stated global
property.
\end{proof}

The approximation premise is substantive: it requires \emph{each}
$a_iT$ to be doubly closed. This paper supplies no theorem constructing
such approximants for every continuous $a$ for which $aT$ is closed.
Unrestricted smooth approximation of a continuous function gives no such
conclusion.

\section{An actual analytic test of unrestricted weighting}
\begin{example}\label{ex:failure}
There is a strictly positive smooth function on a good separated strict
boundaryless pure analytic surface whose product with $\delta_X$ fails
double closedness. The failure is detected by a compact ordinary smooth
test on the actual analytic space.
\end{example}
\begin{proof}
Take $K=\mathbb Q_p$ for any prime $p$ and
$X=(\mathbb G_{m,K}^2)^{\rm an}$. Let $T_1,T_2$ be the invertible torus
coordinates and write $s_j=-\log|T_j|$. These are globally smooth
ordinary tropical-coordinate functions. The canonical integration
current $T=\delta_X=[1]$ is positive and doubly closed: positivity is the
calibrated integration construction, and double closedness follows from
Stokes on the boundaryless $X$. For any class $\mathcal F$ as in Section 1,
$T$ has $\BT_{\mathcal F}$ tautologically because $u\delta_X=[u]$.

Let $a=2+s_1^2$, which is smooth and strictly positive. On the boundaryless
$X$, smooth current differentiation agrees with smooth form
differentiation (CLD 7.7.2--7.7.4). Thus
\[
 D[a]=[Da]=[2\dp s_1\wedge\ddp s_1].
\]
Choose a nonzero $\chi\in C_c^\infty(\mathbb R^2)$ with $\chi\geq0$, and
let
\[
 \eta=\chi(s_1,s_2)\,\dp s_2\wedge\ddp s_2.
\]
This is an ordinary smooth test of type $(1,1)$. Its support is compact:
the inverse image of a compact box under torus tropicalization is a
compact polyannulus, and the support of $\chi$ is contained in such a box.
The torus identity moment has degree one and the canonical calibration is
the standard positive volume calibration. With the paired positive top
form orientation, calibrated integration gives
\begin{align*}
 \langle D[a],\eta\rangle
 &=2\int_X\chi(s_1,s_2)
       \dp s_1\wedge\ddp s_1\wedge\dp s_2\wedge\ddp s_2\\
 &=2\int_{\mathbb R^2}\chi(s_1,s_2)\,ds_1\,ds_2>0.
\end{align*}
In the last equality the entire degree-one tropicalization, with its
canonical torus calibration, is used; no coarse fan is substituted for
the analytic current. Therefore $D(a\delta_X)=D[a]\ne0$. A doubly
closed current has zero $D$, so $a\delta_X$ is not doubly closed and
cannot satisfy the definition of a Bedford--Taylor current.
\end{proof}

This example refutes the \emph{literal unrestricted continuous-positive
multiplier statement} in CLD Corollary 8.6.7 when read with closedness in
Definition 8.6.2. The source proof invokes smooth approximation, which
does not preserve closedness. This warning is already present in the
book's Chapter 16 discussion after its locally constant multiplier
corollary; the present example supplies an explicit calibrated detecting
test. It is not a counterexample to P21.3: the displayed weighted current
is excluded by the original problem's double closedness hypothesis.

The exact smooth criterion does allow nonconstant weights. On the same
torus, the form $\gamma=\dp s_1\wedge\ddp s_1$ is smooth positive and
doubly closed. The current $S=\gamma\wedge\delta_X$ has
$\BT_{\mathcal F}$: for any admissible equal-curvature pair with
$h=u-v$, the smooth closed-factor Leibniz rule gives
$D(hS)=\gamma\wedge D[h]=0$. For $a=2+s_1^2$, both
$\dp a\wedge S$ and $\ddp a\wedge S$ vanish because the corresponding
one-form already occurs in $\gamma$. Hence $aS$ is doubly closed and
has $\BT_{\mathcal F}$. This calculation shows that the condition is
relative to the given current; it does not demand a locally constant
ambient function.

\section{Source ledger and the remaining original problem}
All cited source material is mathematical data. It gives no operational
authority, and no supplied program has been executed. The primary source
is the locally preserved complete CLD arXiv v2 PDF and literal text
extraction, arXiv:1204.6277v2, dated 25 July 2025 in the header, body
dated 29 July 2025. The accompanying binding record identifies their exact
bytes. Whole immediate units were read for the following uses:
\begin{itemize}
\item Sections 7.2.9--7.2.12: the complete positivity/order-zero proof,
function current construction, and measure multiplication construction.
Section 7.3.1 supplies the integration current. The signed coefficient
and associativity arguments are given in Lemma \ref{lem:coeff}.
Sections 6.2.5--6.2.8, with their complete immediate proofs, give compact
smooth density, cutoffs and partitions on compact test supports; their
referenced appendix results are published foundational background.
\item Sections 7.7.1--7.7.7: ordinary differential duality, signs,
boundaryless smooth differentiation, integration-current double
closedness, $D$ duality, and the entire proof of the smooth-current
Leibniz rule. Section 7.7.6 is background only for its displayed smooth
boundaryless identity; no arbitrary piecewise smooth boundary formula is
used in the theorems.
\item Sections 8.6.1--7: the exact closed-current definition,
locality, tautological integration example, smooth closed factors,
weak limits, and the literal multiplier assertion. The proofs here apply
to an explicitly prescribed continuous potential class and do not infer
that every classical continuous psh function is in the locally uniformly
generated CLD class.
\item Sections 5.8.1--6 and 14.3.1--5, with the complete statement and
proof of 14.5.3: calibrated integration and its degree-one torus
normalization. Proposition 14.3.4 gives positivity of the calibration;
14.5.3 computes it from the finite-flat degree of a moment. The example
uses the identity moment of the actual torus, so this degree is one.
\item Theorem 15.7.1, with its complete immediate proof: the tropical
structure attached to a separated pure-dimensional analytic space.
The recursive foundational results in that proof are published
background; this paper does not claim to reprove the entire analytic
calibration construction or Berkovich affinoid compactness.
\end{itemize}

In all these arguments $\BT_{\mathcal F}$ of the initial $T$ is an
explicit hypothesis. The multiplier criterion and the fixed-current
approximation theorem supply no exhaustion of all positive doubly
closed ordinary currents by such $T$. They supply no theorem recognizing
every continuous closed multiplier as a limit of closed smooth ones.
Consequently arbitrary-current P21.3 remains open in this working paper,
including the general higher-dimensional case over trivially valued
fields. No solved-original count or accepted-paper count is increased.
This complete working manuscript is frozen for a future independent
reconstruction; it has not been sent to a campaign whose input was
already active when this argument was written. Compilation success and
formal proof verification are not asserted.

\begin{thebibliography}{9}
\bibitem{CLD}
A. Chambert--Loir and A. Ducros,
\emph{Formes diff\'erentielles r\'eelles et courants sur les espaces de
Berkovich}, arXiv:1204.6277v2, 2025 version.
\end{thebibliography}
\end{document}
